\chapter{Expériences aléatoires et loi binomiale}
\section*{Exercice 21.1 -- Reconnaître le modèle}
[Critique] On tire 5 cartes sans remise d'un jeu de 32 cartes. Le nombre d'as suit-il une loi binomiale ?

\section*{Exercice 21.2 -- Probabilité exacte}
[Base] Si $X\sim\mathcal B(5,1/3)$, calculer $P(X=2)$ et $P(X\ge1)$.

\section*{Exercice 21.3 -- Atteindre une fiabilité}
[Méthode] Des essais indépendants réussissent avec probabilité $0{,}2$. Combien en faut-il pour avoir au moins une réussite avec probabilité au moins $0{,}95$ ?

\section*{Exercice 21.4 -- Intervalle et quantiles discrets}
[Méthode] Pour $X\sim\mathcal B(10,1/2)$, calculer $P(2\le X\le8)$. Trouver le plus petit $b$ tel que $P(X\le b)\ge0{,}95$.

\section*{Exercice 21.5 -- Indépendance par paires}
[Preuve] Deux bits indépendants équiprobables sont tirés. Montrer que les événements « premier bit égal à 1 », « second égal à 1 » et « bits égaux » sont indépendants deux à deux, mais pas mutuellement.

\section*{Exercice 21.6 -- Optimiser un paramètre}
[Raisonnement] Pour $X\sim\mathcal B(n,p)$ et $0<k<n$ fixé, déterminer $p\in[0,1]$ maximisant $P(X=k)$.
